\documentclass{article}

\usepackage[a4paper, margin=2.5cm]{geometry}

\usepackage{graphicx}
\usepackage{amsmath, amssymb}


\usepackage[most]{tcolorbox}
\newtcolorbox{equationbox}{
    colback=gray!10,      % Fundo cinza claro
    colframe=gray!40,     % Borda cinza
    boxrule=0.5pt,        % Espessura da borda
    arc=2mm,              % Cantos arredondados
    left=3mm,
    right=3mm,
    top=2mm,
    bottom=2mm
}
%% \begin{equationbox}\begin{equation} y=x^2 \end{equation}\end{equationbox}


\title{Title of article}
\author{Name of Author}
\date{}



\begin{document}


\maketitle



\section{Title of section}


Write some text here

\begin{equationbox}
\begin{align}
\boldsymbol{\Phi}^{T}\boldsymbol{\Phi}
&\equiv
\mathbf{K}
\end{align}
\end{equationbox}




%\bibliographystyle{plain}
%\bibliography{references}

\end{document}
